\section{Cardinal characteristics and the classical inequalities}
\label{sec:p06-cardinal-characteristics}

An invariant records the least size of a witnessing family, not just the
size of the continuum. The four ideal invariants below ask four different
questions: when small sets stop having a small union, how many cover the
space, how small a nonsmall set can be, and how many small sets dominate
the ideal under inclusion. Their quantifier order is part of the definition.
This development uses the separately identified author sources
\cite[\S522, version of 11 July 2023]{Fremlin522Author2023} and
\cite[\S2, arXiv version 1]{BlassReductions1994V1}.
It is not a replacement for the complete Blass and Bartoszy\'nski
\emph{Handbook of Set Theory} chapters.
The separately frozen prepublication survey
\cite[\S\S2--3, arXiv version 1]{BartoszynskiInvariants1999V1}
also proves definable morphism statements stronger than the cardinal
comparisons alone; those stronger conclusions are not assumed here.

\subsection{Ideals, witnesses, and eventual domination}

\begin{hbconvention}{def:p06-cc-standing}{Standing domains and smallness}
Work in \(\mathsf{ZFC}\). An ideal \(I\subseteq\mathcal P(X)\) is nonempty, proper,
downward closed and closed under finite unions; it contains every singleton.
A \(\sigma\)-ideal is also closed under countable unions. Write
\(\mathfrak c=|\mathbb R|\). Here \(\mathcal N\) is the
ideal of subsets of \(\mathbb R\) of Lebesgue outer measure zero, and
\(\mathcal M\) is the ideal of meager subsets of
\(\mathbb R\). The source's \(\mathcal N,\mathcal M\) are typed here as
ideals, not as structures or ground models. No measurability or Baire
property is required of an arbitrary set whose membership in either ideal
is being tested.
The measure inputs are the usual Lebesgue outer measure, countable
subadditivity, lengths of intervals, open covers of null sets, and the
finite binary-cylinder probabilities. A complete construction of Lebesgue
measure is not claimed as part of this characteristic packet. The needed
complete-metric Baire argument is supplied in the annex.
\end{hbconvention}

\begin{hbdefinition}{def:p06-cc-ideal-four}{The four ideal invariants}
For the standing ideal \(I\), define
\[
\begin{aligned}
\operatorname{add}(I)&=\min\{|\mathcal A|:\mathcal A\subseteq I,
                       \ \bigcup\mathcal A\notin I\},\\
\operatorname{cov}(I)&=\min\{|\mathcal A|:\mathcal A\subseteq I,
                       \ \bigcup\mathcal A=X\},\\
\operatorname{non}(I)&=\min\{|A|:A\subseteq X,\ A\notin I\},\\
\operatorname{cof}(I)&=\min\{|\mathcal B|:\mathcal B\subseteq I,
                       \ (\forall A\in I)(\exists B\in\mathcal B)\ A\subseteq B\}.
\end{aligned}
\]
The family of all singleton subsets ensures that the first two minima exist; \(X\) witnesses
the third and \(I\) itself the fourth. A family satisfying the last
condition is a cofinal base, not a cover alone.
Fremlin writes \(\operatorname{cf}(I)\) for this cofinal-base minimum;
the notation \(\operatorname{cof}(I)\) used in the other identified
presentation denotes the same minimum for the same ideal under inclusion.
This ideal-valued use of \(\operatorname{cf}\) is distinct in type from
ordinal cofinality \(\operatorname{cf}(\alpha)\).
\end{hbdefinition}

\begin{hbtheorem}{thm:p06-cc-ideal-bounds}{proved=annex:p06.cc.ideal-bounds}{The general ideal square and regular additivity}
For the standing ideal,
\[
\operatorname{add}(I)\leq\operatorname{cov}(I)\leq\operatorname{cof}(I),
\qquad
\operatorname{add}(I)\leq\operatorname{non}(I)\leq\operatorname{cof}(I).
\]
For a \(\sigma\)-ideal, \(\aleph_1\leq\operatorname{add}(I)\), and
\(\operatorname{add}(I)\) is a regular cardinal. A Borel cofinal base on
\(\mathbb R\) gives \(\operatorname{cof}(I)\leq\mathfrak c\).
These assertions do not say that every covering or uniformity number is
regular.
\end{hbtheorem}

\begin{hbdefinition}{def:p06-cc-bdsr}{Bounding, domination, splitting, and reaping}
For \(f,g\in\omega^\omega\), write \(f\leq^*g\) when
\(f(n)\leq g(n)\) for all sufficiently large \(n\). Define
\[
\mathfrak b=\min\{|F|:F\subseteq\omega^\omega\text{ has no }\leq^*
  \text{-upper bound}\},\qquad
\mathfrak d=\min\{|D|:(\forall f\in\omega^\omega)(\exists g\in D)f\leq^*g\}.
\]
A set \(S\subseteq\omega\) splits an infinite \(A\subseteq\omega\)
if both \(A\cap S\) and \(A\setminus S\) are infinite.
The splitting number \(\mathfrak s\) is the least size of a family that
splits every infinite set. The reaping number \(\mathfrak r\) is the
least size of a family of infinite sets that no one set splits simultaneously.
Blass's source calls the latter the \emph{unsplitting} number as well.
\end{hbdefinition}

\begin{hbexample}{def:p06-cc-eventual-failure}{Two quantifier distinctions}
The functions \(f_k\) with \(f_k(0)=k\) and \(f_k(n)=0\) for
\(n>0\) are all eventually bounded by zero, but have no pointwise
bound. A countable diagonal bound must therefore be stated with
\(\leq^*\), not with \(\leq\) on every coordinate.
Also, the family of all functions zero at every odd coordinate is
unbounded: against a proposed \(g\), use \(f(2n)=g(2n)+1\).
It is not dominating, because no member eventually dominates the
constant function one. This distinguishes the two witnessing properties;
it does not prove \(\mathfrak b<\mathfrak d\).
\end{hbexample}

\begin{hbtheorem}{thm:p06-cc-combinatorial}{proved=annex:p06.cc.combinatorial}{Combinatorial witness bounds}
\[
\aleph_1\leq\mathfrak b\leq\mathfrak d\leq\mathfrak c,
\qquad
\aleph_1\leq\mathfrak s\leq\mathfrak d,
\qquad
\mathfrak b\leq\mathfrak r\leq\mathfrak c.
\]
The cardinal \(\mathfrak b\) is regular. The interval-partition construction
in the proof simultaneously supplies the splitting witnesses and the
dual reaping obstruction; it does not assert a comparison between
\(\mathfrak s\) and \(\mathfrak r\).
\end{hbtheorem}

\begin{hbtheorem}{thm:p06-cc-real-ideals}{proved=annex:p06.cc.real-ideals}{The real ideals and their bases}
Both standing real ideals are proper \(\sigma\)-ideals containing all
singletons and having Borel cofinal bases. Their four invariants therefore
lie between \(\aleph_1\) and \(\mathfrak c\).
There is a dense \(G_\delta\) null set \(E\) whose complement is meager
and conull. Consequently neither of these two ideals includes the other.
\end{hbtheorem}

\subsection{The category and translation constructions}

\begin{hbtheorem}{thm:p06-cc-category-domination}{proved=annex:p06.cc.category-domination}{Small bounded families and meager covers}
\[
\mathfrak b\leq\operatorname{non}(\mathcal M),
\qquad
\operatorname{cov}(\mathcal M)\leq\mathfrak d.
\]
Eventually bounded cones give the Baire-space argument. On the real line,
inverse-distance codes from the rationals give direct closed nowhere dense
witnesses; the countable rational set is handled separately.
\end{hbtheorem}

\begin{hbtheorem}{thm:p06-cc-category-coding}{proved=annex:p06.cc.category-coding}{Coding a bound from a meager container}
There are maps \(\Phi:\omega^\omega\longrightarrow
\mathcal M\) and
\(\Psi:\mathcal M\longrightarrow\omega^\omega\) such that
\[
 \Phi(f)\subseteq A\quad\Longrightarrow\quad f\leq^*\Psi(A).
\]
In particular
\(\operatorname{add}(\mathcal M)\leq\mathfrak b\)
and \(\mathfrak d\leq\operatorname{cof}(\mathcal M)\).
The closed witness \(\Phi(f)\) is essential: the proof uses Baire's theorem
inside that witness, not an unjustified bound on all meager sets of functions.
\end{hbtheorem}

\begin{hbtheorem}{thm:p06-cc-category-identities}{proved=annex:p06.cc.category-identities}{The two category identities}
\[
\operatorname{add}(\mathcal M)
 =\min\{\mathfrak b,\operatorname{cov}(\mathcal M)\},
\quad
\operatorname{cof}(\mathcal M)
 =\max\{\mathfrak d,\operatorname{non}(\mathcal M)\}.
\]
The reverse inequalities are witnessed by a common translation avoiding
the rational shifts of closed meager covers, followed by one eventual
bound or a dominating family.
\end{hbtheorem}

\begin{hbtheorem}{thm:p06-cc-cross-covers}{proved=annex:p06.cc.cross-covers}{The translation duality}
\[
\operatorname{cov}(\mathcal N)
 \leq\operatorname{non}(\mathcal M),\qquad
\operatorname{cov}(\mathcal M)
 \leq\operatorname{non}(\mathcal N).
\]
The covering sets are translates of the null comeager set or its meager
conull complement. The proof applies to arbitrary nonmeager or nonnull
witnesses, without assuming their regularity.
\end{hbtheorem}

\subsection{Localization and the remaining diagram edges}

\begin{hbdefinition}{def:p06-cc-slaloms}{Localization by a finite-width slalom}
A slalom is a function \(S\) with \(S(n)\subseteq\omega\) and
\(|S(n)|\leq 2^n\). It localizes \(f\in\omega^\omega\), written
\(f\in^*S\), if \(f(n)\in S(n)\) eventually. An empty coordinate is
allowed; the event ``eventually'' cannot be replaced by ``infinitely often.''
\end{hbdefinition}

\begin{hbexample}{def:p06-cc-localization-failure}{Why infinitely many coordinates do not suffice}
In the interval-code construction below, a valid code for \(\{0\}\)
can use small intervals around zero at odd coordinates and small
intervals inside \((1,2)\) at even coordinates. The point zero is in
the limsup of this code. A slalom taking only the coded even values,
and empty at every odd coordinate, matches the code infinitely often;
its own interval limsup misses zero. Thus infinitely-often matching
cannot justify the containment conclusion of the null-code map.
\end{hbexample}

\begin{hblemma}{thm:p06-cc-null-localization}{proved=annex:p06.cc.null-localization}{Two directions of null-set localization}
There are maps \(a:\omega^\omega\to\mathcal N\),
\(b:\mathcal N\to\{\text{slaloms}\}\),
\(c:\mathcal N\to\omega^\omega\), and
\(e:\{\text{slaloms}\}\to\mathcal N\) such that
\[
 a(f)\subseteq A\ \Longrightarrow\ f\in^*b(A),\qquad
 c(A)\in^*S\ \Longrightarrow\ A\subseteq e(S).
\]
The first direction uses independent small open events and a positive
compact set outside \(A\); the second uses summable finite rational-interval
covers. Neither map identifies null sets with meager sets.
\end{hblemma}

\begin{hbtheorem}{thm:p06-cc-meager-null-transport}{proved=annex:p06.cc.meager-null-transport}{Transporting meager containment through null containment}
There are maps \(u:\mathcal M\to
\mathcal N\) and
\(v:\mathcal N\to\mathcal M\) satisfying
\[
 u(A)\subseteq B\quad\Longrightarrow\quad A\subseteq v(B).
\]
Therefore
\(\operatorname{add}(\mathcal N)\leq
  \operatorname{add}(\mathcal M)\) and
\(\operatorname{cof}(\mathcal M)\leq
  \operatorname{cof}(\mathcal N)\).
This is the localization argument behind the two Bartoszy\'nski--
Raisonnier--Stern inequalities in Fremlin's presentation.
\end{hbtheorem}

\begin{hbtheorem}{thm:p06-cc-cichon}{proved=annex:p06.cc.cichon}{The full classical inequality network}
The following generating inequalities, together
with their transitive consequences, give the classical Cicho\'n diagram:
\[
\begin{gathered}
\aleph_1\leq\operatorname{add}(\mathcal N),\\
\operatorname{add}(\mathcal N)\leq\operatorname{add}(\mathcal M)
 \leq\mathfrak b\leq\mathfrak d,\\
\mathfrak d\leq\operatorname{cof}(\mathcal M)
 \leq\operatorname{cof}(\mathcal N)\leq\mathfrak c,\\
\operatorname{add}(\mathcal N)\leq\operatorname{cov}(\mathcal N)
 \leq\operatorname{non}(\mathcal M)\leq\operatorname{cof}(\mathcal M),\\
\operatorname{add}(\mathcal M)\leq\operatorname{cov}(\mathcal M)
 \leq\operatorname{non}(\mathcal N)\leq\operatorname{cof}(\mathcal N),\\
\mathfrak b\leq\operatorname{non}(\mathcal M),\qquad
\operatorname{cov}(\mathcal M)\leq\mathfrak d.
\end{gathered}
\]
The two additional category identities are those of
\cref{thm:p06-cc-category-identities}. The diagram is a network of
inequalities, not a chain asserting that all displayed numbers are equal.
\end{hbtheorem}

\begin{hbcorollary}{thm:p06-cc-ch-collapse}{proved=annex:p06.cc.ch-collapse}{What CH does and does not settle}
Under \(\mathsf{CH}\), all ten classical diagram characteristics, and
\(\mathfrak s,\mathfrak r\), equal \(\aleph_1\).
\end{hbcorollary}

\begin{proofboundary}{def:p06-cc-consistency-boundary}{Absolute inequalities are not separation proofs}
The proofs here are \(\mathsf{ZFC}\) witness constructions. They do not
assert that the maps chosen on arbitrary null or meager sets are Borel
or continuous maps on a coding space. Bartoszy\'nski's 1999 presentation
uses Borel master sets, parameter-uniform covers and selection arguments
to obtain that stronger definable version of the diagram. Mere choices
of a closed cover or a compact set in this packet do not discharge those
additional duties.
The proofs here also do not
establish any strict separation of characteristics, any completeness claim
for the list of consistent assignments, or a characteristic-specific
iteration theorem. Forcing \(\mathfrak c=\kappa\) with Cohen conditions
does not by itself compute \(\operatorname{add},\operatorname{cov},
\operatorname{non},\operatorname{cof},\mathfrak b,\mathfrak d\): each value
requires its own preservation and witness arguments. The complete allocated
source treatments, their variants and their consistency constructions
remain separate obligations.
\end{proofboundary}
